\slajdid{09_3-s063}
\begin{frame}[label=09_3-s063]
\gusttekst\setlength{\abovedisplayskip}{3pt}\setlength{\belowdisplayskip}{3pt}\setlength{\abovedisplayshortskip}{2pt}\setlength{\belowdisplayshortskip}{2pt}Slučaj $V_C=+15V,V_I=0V$\begin{center}\input{slike/tikz/s063_bilateralni_prekidac.tex}\end{center}Da vidimo za N tranzistor. Ima sigurno uslove za provođenje. Ako mu je sors na strani 1 struja kroz otpornik $R_L$ bi tekla iz mase pa bi napon na drejnu tranzistora bio negativan, što nije moguće. Ako pretpostavimo da mu je sors na strani 2, struja kroz otpornik $R_L$ bi tekla u masu, pa bi napon na sorsu bio veći od napona na drejnu što takođe nije moguće. Nebitno na kojoj strani je sors, očigledno je da N tranzistor mora da radi sa strujom koja je jednaka nuli odnosno da radi u omskoj oblasti sa malom dinamičkom otpornosti. Potpuno identično razmišljanje možemo da izvedemo i za P tranzistor, odnosno dolazimo do zaključka da će $V_O=0V$.
\end{frame}
